\section{Integral transfer to a tame coefficient-prime-power branch}
\label{sec:cubic-hsieh-tame-measure}

\begin{proposition}[AC-G67-measure: the common Eisenstein measure preserves the congruence]
\label{AC-G67-measure}
Keep the imaginary quadratic split-$p$ setting and all retained
published inputs of AC-K29-Hsieh-aux.  Choose its extra prime $q$
outside the original conductor and with $q\equiv1\pmod{p^r}$,
where $r\ge1$.  Let $\vartheta$ be a rational Dirichlet character
of exact order $p^r$ and conductor $q$, restricted to $G_K$.
Write $\beta'=\beta\vartheta$.

There is a common sufficiently deep prime-to-$p$ level $n$ at
which the integral ordinary Eisenstein series for $\beta$ and
$\beta'$ lie in the same ordinary weight module after coefficient
extension and satisfy
\[
        E^{\mathrm{ord}}(\beta'\varepsilon_\Lambda\mid1,n)
        \equiv E^{\mathrm{ord}}(\beta\varepsilon_\Lambda\mid1,n)
                       \pmod{\mathfrak m_{\mathcal O'}}.       \tag{G67.1}
\]
The first series is primitive at the maximal ideal of
$\Lambda'=\mathcal O'[[\Gamma_K]]$, and all its boundary
constant terms lie in that maximal ideal.  Hence the same
fundamental exact sequence and subtraction argument as in
Hsieh Theorem~7.18 give nonzero ordinary cuspidal Hecke support
for the branch $\beta'$ at this common level.

This statement uses the integral measure in Hsieh
Proposition~5.14 and Corollary~5.15.  It does not assert the
general removal of the hypothesis $p\nmid\#\Delta$ from
Burungale--He--Tian--Ye Theorem~6.5, nor does it project the
cuspidal support to a specified $q$-inertial type.
\end{proposition}
\begin{proof}
\proofstep{Choose the common level and coefficient ring.}
The cyclic group $(\mathbb Z/q\mathbb Z)^\times$ has a quotient
of order $p^r$, giving the stated character.  Its restriction
to $G_K$ still has that order because $[K:\mathbb Q]=2$ is
prime to $p$.  It is unramified above $p$, and every value is
congruent to $1$ in a coefficient discrete valuation ring
$\mathcal O'$ containing its values.  Its conductor over $K$
is $q\mathcal O_K$, since $q$ splits and its local character
is nontrivial at each of the two primes.

Let $c$ be the prime-to-$p$ conductor of $\beta$ and put
$c'=c q\mathcal O_K$.  In Hsieh (a1), take the split-prime
product for the original branch to contain $q^2$; this is
still a product of split primes.  Then a level of the form
\[
 n=(s q^2c\bar cD_{K/\mathbb Q})^{m_0}
   =(s c'\bar c'D_{K/\mathbb Q})^{m_0}
\]
has the required form for both conductor bounds.  Deepen the
level so that (a2)--(a3) and the original Fourier--Jacobi unit
construction hold, as permitted in the proof of Proposition~7.13
and AC-K29-Hsieh-aux.  The original Fourier--Jacobi triple is
chosen with $q$ outside its fixed bad primes.  No new prime
outside the chosen support is introduced by this common-depth
choice.

\proofstep{Increasing the conductor bound leaves the original normalized series unchanged.}
Here both constructions use this same $n$, so that their set
$S=\{v:v\mid n\}$ already contains $q$.  The classical local
sections in Hsieh \S5.3 depend on $\chi$, $n$ and $S$; the
conductor bound $c$ only specifies which characters may be
used.  In particular the section at $q$ is already the
ramified-level section (5.4)--(5.5), even when $\chi$ itself is
unramified there.  The normalizing factor preceding (5.10) is
$\Lambda_2^{S\cup S_p}(\chi,s)$ in both constructions.
Thus enlarging $c$ to $c'$ does not remove a second Euler
factor or multiply this fixed-level series by any new factor.

This also proves equality of the integral measures without
relying only on their notation.  Let $\varpi:G_{c'}\to G_c$
be restriction.  The pushforward under $\varpi$ of the
$c'$-measure in Theorem~5.8 has, at every character in
$X^+_{G_c}$, the same interpolation value $\widehat E(\chi,n)$
as the $c$-measure.  The uniqueness assertion in that theorem
therefore identifies the measures.  The ordinary pullback in
the proof of Proposition~5.14 uses the same finite set
$X_W(n)$, the same level $K_P(n)$ and the same ordinary
projector; it consequently preserves this identification.
Corollary~5.15 then identifies the $\beta$ specialization of
the common $c'$-measure with the original $\beta$ branch at
this fixed $n$, exactly rather than up to a scalar.

\proofstep{Use a single integral measure before specializing the finite branch.}
For the conductor bound $c'$ and this fixed level $n$, Hsieh
Proposition~5.14 gives an integral ordinary Eisenstein measure
on $G_{c'}=\operatorname{Gal}(K_{c'p^\infty}/K)$.
Corollary~5.15 forms the branch series by integrating
\[
        g\longmapsto\beta(g)\,
                  \varphi(g|_{K_\infty})
\]
against this measure, for continuous integral test functions
$\varphi$ on $\Gamma_K$.
For the other branch the integrand is multiplied by
$\vartheta(g)$, which is congruent to $1$ modulo
$\mathfrak m_{\mathcal O'}$ for every $g$.
Integration against an integral measure preserves this
congruence: this is immediate first for locally constant
functions, which are finite sums over compact open sets, and
then for continuous functions by uniform approximation and
completeness.  Thus all the integral coefficients of the two
measure-valued forms are congruent, proving (G67.1).

Their ordinary weight modules are the same.  Indeed the weight
homomorphism in Hsieh (5.25) uses the two local restrictions
above $p$, and $\vartheta$ is unramified there.  In particular
it is trivial on the local unit groups that define that
homomorphism.  We have only extended the coefficient ring to
include the finite character values.

\proofstep{Transfer the unit and the nonunit boundary.}
The original integral Fourier--Jacobi functional
$\ell_{a,m,\kappa}$ is a unit on the original branch at the
chosen level.  It is defined on the common ordinary module
and is integral after coefficient extension.  Applying it to
(G67.1) shows that it is also a unit on the $\beta'$ branch.
This establishes primitivity without applying an unproved
generalization of the repaired nonvanishing theorem to a
finite extension of degree divisible by $p$.

By AC-K29-Hsieh-aux, every boundary constant term of the original
series is in the maximal ideal of $\Lambda'$: Hsieh
Proposition~6.6 identifies their ideal with the product of the
two imprimitive $p$-adic $L$-functions, and the added $q$-Euler
factor makes that product a nonunit.  The boundary map is
integral and linear, so (G67.1) puts every boundary constant
term of the $\beta'$ series in the same maximal ideal.
Proposition~6.6 for the latter series has a distinguished
constant term equal to its product $L_D^S L_K^S$ times a unit;
hence that product is a nonunit as well.

Now repeat the complete proof of Hsieh Theorem~7.18 after its
primitivity step.  The fundamental exact sequence lifts the
boundary divided by this product to an ordinary integral form
$E_0$.  Subtracting $E_0L_D^SL_K^S$ produces a cusp form on
which the same Fourier--Jacobi functional is still a unit.
The resulting Hecke quotient maps onto
$\Lambda'/(L_D^SL_K^S)\ne0$, proving nonzero cuspidal support.
No averaging over the group of order $p^r$ is used at any point.
\end{proof}

\paragraph{Exact local-type and geometric scope.}
At $q$, the three desired finite inertia characters for this
branch are $\vartheta^{-1},1,\vartheta$, up to the fixed
Weyl ordering.  They are distinct in characteristic zero and
reduce to $1$.  The proposition establishes support in the
whole common-level module.  It does not establish support
inside that exact character-type module, nor boundary
surjectivity there; a rational projector with denominator
$p^r$ does not supply those statements integrally.
The separate symmetric Hodge-order obstruction of
AC-G64-weight is also retained.
The finite character is a rank-one lisse system on the open
arithmetic curve after deleting $pq$, in the previously audited
SGA~1, Expos\'e~V, \S\S4 and~7 framework.  This dictionary
does not assert an unverified descent across the boundary of
the unitary Shimura variety.
